% =============================================================================
% Primer Template (Conceptual, Non-Problem Content)
% =============================================================================
% Usage:
% - Copy into: latex/sections/<domain>/primer.tex (or a more specific name)
% - Keep code in external files under: problems/<domain>/snippets/<slug>/
% - From primer.tex, reference code with paths like:
%     ../../../problems/<domain>/snippets/<slug>/solution.cpp
%     ../../../problems/<domain>/snippets/<slug>/solution.py
% - Keep this focused on recognition, invariants, core ops, complexity, pitfalls
.
% =============================================================================

% Replace <Topic> with the data structure / concept name
\subsection{Primer — <Topic>}

\paragraph{What It Is}
% Concisely describe the structure/concept, its representation, and when it is u
sed.
% Emphasise the mental model (e.g., nodes + pointers; contiguous array + indices
).

\paragraph{Recognition Patterns}
% How to identify that this structure is implied by a question. Provide 3–6 cues
.
\begin{itemize}[nosep]
  \item <Cue 1>
  \item <Cue 2>
  \item <Cue 3>
\end{itemize}

\paragraph{Core Operations}
% Briefly define the main operations and their expected complexity.
% Include short pseudocode blocks when clarity benefits.
\begin{CodeBlock}[style=pseudo,caption={<Topic>: Core Operations (pseudocode)}]{
18}
<op1>(...):
  <pseudocode>

<op2>(...):
  <pseudocode>
\end{CodeBlock}

\paragraph{Invariants \& Correctness}
% State the key invariants maintained (e.g., head/tail consistency, ordering, sh
ape).
% Sketch why the core operations preserve these invariants.

\paragraph{Complexity}
% Summarise asymptotic time/space for each core operation.
\begin{itemize}[nosep]
  \item <op1>: time $O(\cdot)$, space $O(\cdot)$
  \item <op2>: time $O(\cdot)$, space $O(\cdot)$
\end{itemize}

\paragraph{Common Pitfalls}
% Practical footguns and frequent mistakes.
\begin{itemize}[nosep]
  \item <Pitfall 1>
  \item <Pitfall 2>
  \item <Pitfall 3>
\end{itemize}

% -----------------------------------------------------------------------------
% Example Implementations (external files; adjust paths for your primer location
)
% -----------------------------------------------------------------------------
% From latex/sections/<domain>/primer.tex, use ../../../ to reach repo root.

% \bigskip
% \noindent\textbf{C++ Example}
% \lstinputlisting[
%   style=cpp,
%   caption={<Topic> — Minimal C++ Example}
% ]{../../../problems/<domain>/snippets/<slug>/solution.cpp}

% \bigskip
% \noindent\textbf{Python Example}
% \lstinputlisting[
%   style=python,
%   caption={<Topic> — Minimal Python Example}
% ]{../../../problems/<domain>/snippets/<slug>/solution.py}
